\section{Evaluation by Flaw Injection}
\label{sec:evaluation}

\subsection{Local ground truth without the counterfactual}
The defining difficulty in evaluating a causal-identification auditor is that the
quantity of ultimate interest (the true causal effect of a carbon price, an
emissions cap, or an emission standard) is never observed. We therefore do not
attempt to certify that \sys{} recovers true effects. Instead we construct
\emph{local} ground truth by controlled perturbation. Beginning from a credible
policy-evaluation paper, we inject exactly one known flaw that targets a single
identification dimension, drawn from the flaw taxonomy
(\texttt{flaw\_taxonomy.yaml}). Because we control which dimension we damaged, we
know the correct audit outcome for that dimension, without ever needing the true
effect. The bottom band of Figure~\ref{fig:arch} shows the loop: the clean and
injected versions are each audited, and three metrics are computed.

\begin{itemize}
  \item \textbf{Detection}: does the injected version raise risk on the targeted
        dimension above a threshold?
  \item \textbf{False alarm}: does the clean version stay quiet on that dimension?
  \item \textbf{Localization}: is the targeted dimension the one the auditor flags
        most strongly?
\end{itemize}

\subsection{Realistic perturbation, not sentinels}
How a flaw is injected is a methodological choice that matters. A na\"ive injector
appends a sentence that \emph{names} the flaw: e.g.\ ``treated and control have
non-zero pre-period coefficients with significant leads.'' If the detector then
searches for exactly those phrases, detection becomes a tautology: the injector
plants the very string the detector greps for. In our initial implementation this
produced a perfect-looking but meaningless detection rate of $1.000$.

We replace this with \emph{realistic structural perturbation} of two kinds.
\emph{Omission} flaws \emph{remove} the section, figure, or table that supplies the
support for a dimension (e.g.\ deleting the placebo-robustness section or the
balance table), modelling a paper that simply omits the check; an auditor can flag
this non-circularly by noticing the \emph{absence} of evidence. \emph{Commission}
flaws \emph{rewrite} a section in flawed-but-natural prose, the way a real flawed
paper would describe a mistaken design choice (e.g.\ pooling all province-years in
a two-way fixed-effects regression under staggered policy timing), \emph{without}
containing any of the phrases the detector looks for. We enforce the latter with a
regression test that fails if an injected paper contains the target dimension's
negative-signal phrases, so the circularity cannot silently return. Detection is
thereby a genuine task rather than a string match, which is what makes the results
in \S\ref{sec:results} meaningful.
